% The Tanner graph of Section~\ref{sec:ldpc}, drawn by hand. Input by
% textbook.tex; a sketch and not a rendered figure, so it is outside the
% figure cache and the render cap, which govern generated output only.
\begin{figure}[htbp]
\centering
\begin{tikzpicture}[
    scale=1.0,
    check/.style={fgfactor, minimum size=0.45cm, inner sep=1pt,
                  font=\footnotesize}
]
  \foreach \i in {1,...,6} {
    \node[fgvariable] (c\i) at (1.35*\i, 1.5) {$c_{\i}$};
    \node[fgdata] (l\i) at (1.35*\i, 0.7) {};
    \draw[fglink] (c\i) -- (l\i);
    \node[font=\footnotesize, anchor=north] at (1.35*\i, 0.55) {$L_{\i}$};
  }
  \node[check] (a) at (2.7, 3.1) {$+$};
  \node[check] (b) at (5.4, 3.1) {$+$};
  \node[check] (d) at (8.1, 3.1) {$+$};
  \foreach \bit in {1,2,3,4} { \draw[fglink] (c\bit) -- (a); }
  \foreach \bit in {2,4,5,6} { \draw[fglink] (c\bit) -- (b); }
  \foreach \bit in {1,3,5,6} { \draw[fglink] (c\bit) -- (d); }
  \node[fgnote] at (9.0, 3.1)
    {$\psi_a = \mathbb{1}[\bigoplus_{i \in \partial a} c_i = 0]$};
  \node[fgnote] at (9.0, 0.7)
    {$\psi_i = e^{-c_i L_i}$};
\end{tikzpicture}
\caption{A Tanner graph: bits $c_i \in \{0, 1\}$ as circles, the parity
factors of a row of $\parity$ as squares, and the channel factor of
Equation~\ref{eq:ldpc-llr} as a shaded square carrying the log-likelihood
ratio $L_i$ at each bit. The drawn graph is six bits and three checks of
degree four; the ensemble this repository carries is regular $(3, 6)$, so
every bit sits on three parity factors and every parity factor on six bits.
The graph has cycles --- a random draw of the ensemble does not avoid short
ones --- so Equation~\ref{eq:sum-product} on it is the Bethe approximation
and Equation~\ref{eq:tanh-rule} is that approximation in the log domain. A
sketch of the structure and not a rendering of any instance.}
\label{fig:ldpc:sketch}
\end{figure}
